\documentclass[a4paper,11pt]{article}
\usepackage{amsfonts}
\usepackage{amsmath}
\usepackage{amssymb}
\usepackage{amsthm}
\usepackage{enumerate}
\usepackage[spanish]{keystroke}
\usepackage[spanish]{babel}
\usepackage[T1]{fontenc}
\usepackage[utf8]{inputenc}
\usepackage{lmodern}
\usepackage{pseudocode}
\newtheorem{teo}{Teorema}
\newtheorem{cor}{Corolario}
\newtheorem{prop}{Propiedad}

\title{Laboratorio de Matemática Discreta II \\Algoritmo de Edmonds-Karp}
\author{Fernando Marengo\\
Mariano Mateuci\\
Giovanni Rescia\\
Nahuel Seiler\\
Pablo Ventura}

\begin{document}

\maketitle
%\tableofcontents


\section{Diseño y Modularización}
El diseño esta dividido en cuatro módulos:
\begin{itemize}
  \item main: Para la ejecución principal del programa.
  \item matrizcap: Contiene la implementación de la clase MatrizCap que implementa una nuestra representación de la matriz de capacidades.
  \item api: Contiene la implementación de la clase EstadoEK que implementa Edmonds-Karp.
  \item constants: Con las constantes de uso comun de los demás modulos.
\end{itemize}

\subsection{Módulo main}
Nuestro main utliza la biblioteca optparse para tomar el argumento de la verbosidad al iniciarse, tambien importa el modulo api. Luego crea una instancia de EstadoEK con la verbosidad correspondiente. Luego va llamando a leer\_un\_lado hasta que deje de devolver $1$, para despues empezar a llamar a aumentar\_flujo hasta que deje de devolver $1$. Si devuelve $-1$ imprime un mensaje de error y en otro caso llama a imprimir\_flujo\_maximal.

\subsection{Módulo matrizcap}
Implementa la clase MatrizCap que implementa una nuestra representación de la matriz de capacidades.
\paragraph{Clase MatrizCap}
Utiliza un diccionario para guardar los valores. El diccionario utiliza como clave una tupla con los vertices del lado y como valor una tupla con la capacidad total y el flujo actual.
\subparagraph{Método agregarLado}
Toma como argumentos los vertices $(x,y)$ del lado y su capacidad y lo agrega a la matriz de capacidades. Inicializa el flujo actual en $0$.
Va agregando a los vertices a dos diccionarios que guardan los $\varGamma^+$ y $\varGamma^-$ de cada vertice.
\subparagraph{Método capacidadTotal}
Toma como argumentos los vertices $(x,y)$ de un lado y devuelve su capacidad total.
\subparagraph{Método capacidadActual}
Toma como argumentos los vertices $(x,y)$ de un lado y devuelve su capacidad actual.
\subparagraph{Método flujoActual}
Toma como argumentos los vertices $(x,y)$ de un lado y devuelve su flujo actual.
\subparagraph{Método gammaMas}
Toma como argumento un vertice del network y devuelve una lista con su $\varGamma^+$
Como la lista ya se genero al agregar los lados, esta funcion es de orden constante.
\subparagraph{Método gammaMenos}
Toma como argumento un vertice del network y devuelve una lista con su $\varGamma^-$
Como la lista ya se genero al agregar los lados, esta funcion es de orden constante.
\subparagraph{Método noEstaSaturado}
Toma como argumentos los vertices $(x,y)$ de un lado y devuelve $True$ si el lado no está saturado actualmente.
\subparagraph{Método flujoDesdeFuente}
No toma argumentos. Devuelve el valor de todo el flujo que sale desde la fuente $s$, cuando no hay más caminos aumentantes devuelve el flujo maximal.
\subparagraph{Método flujoForward}
Toma como argumentos el lado $(x, y)$ y el flujo forward que pasa por es lado para modificar la matriz de capacidades sumando ese flujo al lado $(x,y)$.
\subparagraph{Método flujoBackward}
Toma como argumentos el lado backward $(x, y)$ y el flujo forward que pasa por es lado para modificar la matriz de capacidades restando ese flujo al lado $(y,x)$.
\subparagraph{Método \_\_str\_\_}
No toma argumentos. Implementado para devolver una impresión acorde a la especificación al hacer un print de la matriz de capacidades.

\subsection{Modulo api}

\paragraph{Clase EstadoEK}
Implementa el calculo del flujo maximal con Edmonds-Karp.
Utiliza una instancia de MatrizCap, para almacenar las capacidades y los flujos de los lados.
\subparagraph{Método \_\_init\_\_}
Inicializa el objeto. Verifica que la verbosidad sea válida.

\subparagraph{Método leer\_un\_lado}
Toma desde la entrada estandar una expresión respetando la especificación. Utiliza la biblioteca re para parsear mediante expresiones regulares. Cuando la expresión matchea correctamente es agregada a la matriz de capacidades y se retorna 1. En caso contrario se retorna 0.
\subparagraph{Método aumentar\_flujo}
Empieza instanciando un objeto BFS para el estado actual de la matriz de capacidades. Luego pregunta si pudo llegar al resumidero, en ese caso instancia un objeto CaminoAumentante con el bfs actual y la matriz de capacidades, comprueba el invariante de que el largo de cada nuevo camino sea mayor o igual que el anterior. Si no pudo llegar al resumidero calcula el flujo maximal usando flujoDesdeFuente y pide el corte minimal al bfs y luego se lo pasa a la matriz de capacidades para saber su capacidad. En ese momento comprueba el invariante de que la capacidad del corte minimal debe ser igual al flujo maximal. Luego hace las impresiones correspondientes segun la verbosidad.

\subparagraph{Método imprimir\_flujo\_maximal}
Hace la impresion del flujo maximal respetando la especificación y la verbosidad actual.

\subsection{Modulo bfslist}

\paragraph{Clase BFS}

\subparagraph{Método \_\_init\_\_}
Genera el bfs a partir de la matriz de capacidades. Utiliza un diccionario cuyas claves son el vertice y su valor una named tuple con vertice, padre, flujo, esBackward.
Al iterar va generando una lista "iterador" para poder ir agregando elementos durante el mismo bucle evitando el problema de iterar sobre un diccionario que va cambiando de tamaño.

\subparagraph{Método \_flujoBfs}
Toma un vertice, y devuelve el flujo que tiene en el bfs.

\subparagraph{Método devolver}
Toma un vertice, y devuelve el padre que tiene en el bfs.

\subparagraph{Método \_estaEnBfs}
Toma un vertice, y devuelve un booleano que dice si está o no en el bfs.

\subparagraph{resumideroEnBfs}
Devuelve el BfsElem que representa al resumidero, o si no llego None.

\subparagraph{corteMinimal}
Devuelve la lista de los vertices que quedaron sin poder llegar al resumidero.

\subsection{Modulo caminoaumentante}

\paragraph{Clase CaminoAumentante}

\subparagraph{Método \_\_init\_\_}
Calcula el camino aumentante a partir del bfs y la matriz de capacidades.

\subparagraph{Método \_\_len\_\_}
Devuelve el largo del camino para poder ver el invariante del largo de los caminos.

\subparagraph{Método flujo}
Devuelve el flujo del camino aumentante.

\subparagraph{Método \_\_str\_\_}
Devuelve un string con la impresion de la especificacion.

\subsection{Módulo constants}
Incluye una variable con las verbosidades válidas y dos variables con el numero asignado a la fuente y al resumidero.

\section{Ampliacion para tomar entradas alfabeticas}
Implementamos el script $convertidor.py$ que toma desde la entrada estándar una matriz de capacidades expresada de la forma ``xy cap'' donde $x$ e $y$ son los vertices del lados pero con letras como nombres y va devolviendo por la salida estandar lineas donde las letras son reemplazadas por $0$ o $1$ si son son $s$ o $t$ respectivamente, o en otro caso por su código Ascii y separadas por un espacio junto con la capacidad. De esta forma pudimos utilizar test con este formato para la entrada de nuestro programa principal.
\end{document}